% Part: many-valued-logic
% Chapter: syntax-and-semantics
% Section: introduction

\documentclass[../../../include/open-logic-section]{subfiles}

\begin{document}

\olfileid{mvl}{syn}{int}

\olsection{Pambuka}

Ing logika klasik, kita nyinaoni !!{formula}s kang kawangun saka
!!{propositional variable}s kanthi konektif proposisional
$\lnot$, $\land$, $\lor$, $\lif$, lan $\liff$. Nalika netepake
semantik kanggo logika klasik, kita nggunakake rong nilai kabeneran,
$\True$ lan $\False$. Kita napsirake !!{propositional variable}s ing
!!a{valuation}~$\pAssign{v}$, kang masangake nilai kabeneran
$\True$, $\False$ marang !!{propositional variable}s. Saben
!!{valuation} banjur nemtokake nilai kabeneran $\pValue{v}(!A)$
kanggo saben !!{formula}~$!A$; lan !!^a{formula} dicukupi dening
!!a{valuation}~$\pAssign{v}$, $\pSat{v}{!A}$, yen lan mung yen
$\pValue{v}(!A) = \True$.

Logika mawa akeh nilai nggeneralisasi logika klasik kang nduweni
rong nilai kanthi ngidini nilai kabeneran luwih akeh tinimbang
$\True$ lan $\False$. Mula ing logika mawa akeh nilai,
!!a{valuation}~$\pAssign{v}$ yaiku fungsi kang masangake
salah sawijining nilai kabeneran kang diidini marang saben
!!{propositional variable}~$p$. Himpunan nilai kabeneran kang
diidini umume diarani~$V$. Logika klasik kalebu logika mawa akeh
nilai kanthi $V = \{\True, \False\}$; nilai kabeneran
$\pValue{v}(!A)$ diitung nganggo tabel kabeneran kang lumrah
kanggo saben konektif.

Sawise nilai kabeneran liyane ditambahake, ana luwih saka siji
cara kang lumrah kanggo ngitung~$\pValue{v}(!A)$ tumrap konektif
kang kita waca minangka ``lan,'' ``utawa,'' ``ora,'' lan
``yen---mula.'' Dadi logika mawa akeh nilai ora mung ditemtokake
dening himpunan nilai kabeneran, nanging uga dening
\emph{fungsi kabeneran} kang dipilih kanggo saben konektif.
Sawise pilihan iku ditetepake kanggo logika mawa akeh nilai~$\Log L$,
nilai kabeneran $\pValue{v}(!A)[\Log L]$ ditemtokake kanthi unik
dening valuasi, kaya ing logika klasik. Logika mawa akeh nilai,
kaya logika klasik, iku \emph{fungsional tumrap kabeneran}.

Kanthi dhasar semantis iki, kita bisa nerusake kanggo netepake
padanan konsep semantis tautologi, akibat semantis, lan bisa dicukupi.
Ing logika klasik, !!a{formula} iku tautologi yen nilai kabenerane
$\pValue{v}(!A) = \True$ kanggo saben~$\pAssign{v}$. Ing logika
mawa akeh nilai, konsep iki uga kudu digeneralisasi. Sepisanan,
himpunan nilai kabeneran~$V$ ora kudu ngemot~$\True$. Tuladhane,
sawetara logika mawa akeh nilai nggunakake wilangan, kayata kabeh
wilangan rasional ing antarane $0$ lan~$1$, minangka himpunan nilai
kabeneran. Ing kahanan iki, $1$~umume nduweni peran kaya~$\True$.
Ing logika liyane, ora mung siji nanging sawetara nilai kabeneran
nduweni peran iku. Mula saben logika mawa akeh nilai kudu nduweni
himpunan~$V^+$ kang dumadi saka \emph{nilai pinilih}. Kita banjur
bisa kandha yen !!a{formula} dicukupi dening
!!a{valuation}~$\pAssign{v}$, $\pSat{v}{!A}[\Log L]$, yen lan mung yen
$\pValue{v}(!A)[\Log L] \in V^+$. !!^a{formula}~$!A$ iku tautologi
ing logika iki, $\Entails[\Log L] !A$, yen lan mung yen
$\pValue{v}(!A)[\Log L] \in V^+$ kanggo saben~$\pAssign{v}$.
Pungkasan, kita kandha yen $!A$ iku akibat saka himpunan
!!{formula}s, $\Gamma \Entails[\Log L] !A$, yen saben !!{valuation}
kang nyukupi kabeh !!{formula}s ing~$\Gamma$ uga nyukupi~$!A$.

\end{document}
